\setcounter{section}{2}
\section{Strom, Spannung, Widerstand, Leistung, Energie}
\subsection{Strom- und Spannungsmessung II}

\begin{sol}{EI101}
  Siehe Beschreibung am Anfang dieses Kapitels.
\end{sol}

\begin{sol}{EI102}
  Die mit \qq{A}(Ampere) markierte Strommessung muss in den Stromkreis eingeschleift sein. Die mit \qq{V}(Volt) markierte Spannungsmessung muss parallel zum Messobjekt verwendet werden. Dies ist nur bei Antwort (A) der Fall.
\end{sol}

\begin{ohmchapter}

    \begin{description}
        \item[Spannungsmessung] Parallel zum Messobjekt und hochohmig.
        \item[Strommessung] In den Stromkreis eingeschleift und niederohmig.  
    \end{description}

    Zum Merken: Die Spannungsmessung erfolgt parallel zum Messobjekt. Deshalb muss das Messgerät hochohmig sein, sonst würde der Strom nicht durch unsere Schaltung, sondern durch das Messgerät fließen. Bei der Strommessung muss das Messgerät Teil der Schaltung sein, damit der Stromfluss gemessen werden kann. Hier sollte das Messgerät die Schaltung in keiner Weise beeinflussen, deshalb sollte es möglichst wenig Widerstand haben: niederohmig.
  \end{ohmchapter}   

\subsection{Zeigerinstrumente ablesen}



\begin{sol}{EI103}
  Da \SI{10}{\volt} eingestellt sind, verwendest Du am einfachsten die obere Skala, die Werte bis 100 anzeigt.
  Auf der Skala kannst Du den Wert 29 ablesen. Diesen Wert musst Du noch durch 10 teilen, da die Skala ja bis 100 geht.
  Also: $29 / 10 = \SI{2.9}{\volt}$
\end{sol}

\begin{sol}{EI104}
 Wir verwenden die untere Skala, die bis zum Wert 30 anzeigt, und lesen den Wert 8,8 ab. Da das Messinstrument auf \SI{300}{\volt} eingestellt ist und die Skala bis 30 geht, müssen wir mit Faktor 10 multiplizieren.
 $8,8 \cdot 19 = \SI{88}{\volt}$
\end{sol}

\begin{ohmchapter}
Beim Ablesen von analogen Messgeräten stehen oft unterschiedliche Skalen bereit. Du solltest die Skala auswählen, die möglichst gut zu dem Messbereich passt, den Du verwendest, damit Du nur einen Faktor wie z.B. 10/100 etc. berücksichtigen musst.
\end{ohmchapter}   

\subsection{Spitzen- und Effektivwert}

\begin{sol}{EB401}
 Rechnung:
   $$\hat{U} = U_\text{eff} \cdot \sqrt{2} = \SI{230}{\volt} \cdot \sqrt{2} \approx \SI{325}{\volt}$$
\end{sol}

\begin{sol}{EB402}
 Wie in Frage \qref{EB401}, errechnen wir $\hat{U} \approx \SI{325}{\volt}$.

 Rechnung: $$U_\text{SS} = 2 \cdot \hat{U} = 2 \cdot \SI{325}{\volt} \approx  \SI{651}{\volt}$$
\end{sol}


\begin{sol}{EB403}
 Rechnung:
   $$\hat{U} = U_\text{eff} \cdot \sqrt{2} = \SI{12}{\volt} \cdot \sqrt{2} \approx \SI{17}{\volt}$$
   $$U_\text{SS} = 2 \cdot \hat{U} = 2 \cdot \SI{17}{\volt} = \SI{34}{\volt} $$
\end{sol}


\begin{sol}{EB404}
 Rechnung:
   $$U_\text{eff} = \frac{\hat{U}}{\sqrt{2}} = \frac{\SI{12}{\volt}}{\sqrt{2}}  \approx \SI{8.5}{\volt}$$
\end{sol}


\begin{sol}{EB405}
 Rechnung:
   $$U_\text{eff} = \frac{\hat{U}}{\sqrt{2}} = \frac{\SI{1}{\volt}}{\sqrt{2}}  \approx \SI{0.7}{\volt}$$
\end{sol}

\begin{sol}{EB406}
 Du kannst ablesen, dass der Spitze-Spitze-Wert 4 Einteilungen entspricht. Jede Einteilung entspricht \SI{3}{\volt}.
 Somit ist $U_\text{SS} = 4 \cdot \SI{3}{\volt} = \SI{12}{\volt}$
\end{sol}


\begin{sol}{EB407}
    Analog zu Frage \qref{EB406}.
 Du kannst ablesen, dass der Spitze-Spitze-Wert 4 Einteilungen entspricht. Jede Einteilung entspricht \SI{10}{\volt}.
 Somit ist $U_\text{SS} = 4 \cdot \SI{10}{\volt} = \SI{40}{\volt}$
\end{sol}

\begin{ohmchapter}
    In diesem Kapitel geht es um Wechselspannung. Die Spannung ändert sich die ganze Zeit wie eine Sinusfunktion.
    Welche Angaben können wir über solch eine Wechselspannung machen?

    \begin{description}
        \item[Spitzenspannung] Die Spitzenspannung $\hat{U}$ gibt den höchsten Spannungswert der Wechselspannung an.
        \item[Spitze-Spitze-Wert] Der Spitze-Spitze-Wert $U_\text{SS}$ gibt die Spannung von niedrigstem zum höchsten Punkt der Wechselspannung an.        
    \end{description}
    
    Dies sieht wie folgt aus:

    \pic[1.1]{U_SS}
    In der Formelsammlung findest Du im Abschnitt \qq{Wechselspannung} die Formel $U_\text{SS} = 2 \cdot \hat{U}$, d. h., wir verdoppeln die Spitzenspannung einfach, um den Spitze-Spitze-Wert zu erhalten.

    Die Spitzenwerte der Spannung werden allerdings immer nur ganz kurz erreicht und liegen im Kurvenverlauf deutlich darunter. Wie können wir so etwas wie eine Art \qq{Durchschnitt} für die Spannung angeben? 
    Tatsächlich gibt es so etwas Ähnliches wie den Durchschnitt, es ist der sogenannte \textbf{Effektivwert} der Spannung und der wird in Formeln als $U_\text{eff}$ bezeichnet. Wenn mit diesem Effektivwert gerechnet wird, gilt das Ohm’sche Gesetz weiter. Er gibt quasi den Spannungswert an, der gleich viel Energie besitzt wie die Wechselspannung.

    Dies sieht etwas so aus:
    \pic[1.1]{U_EFF}

    In der Formelsammlung findest Du zur Umrechnung die folgende Formel:
    $$\hat{U} = U_\text{eff} \cdot \sqrt{2}$$
    Umgekehrt gilt natürlich: 
    $$U_\text{eff} = \frac{\hat{U}}{\sqrt{2}}. $$

\end{ohmchapter}   

\subsection{Oszilloskop I}

\begin{sol}{EB408}
    Rechnung: 
  $$ f = \frac{1}{T}  = \frac{1}{\SI{0.00005}{\second}} = \SI{20}{\kilo\hertz}$$
\end{sol}


\begin{sol}{EB409}
    Eine Periode hat 4 Einteilungen und eine Einteilung entspricht \SI{3}{\micro\second}. Also ist eine Periode $4\cdot \SI{3}{\micro\second} = \SI{12}{\micro\second}$, also

  $$ f = \frac{1}{T}  = \frac{1}{\SI{0.000012}{\second}} = \SI{83,3}{\kilo\hertz}$$
\end{sol}


\begin{sol}{EB410}
    Eine Periode hat 4 Einteilungen und eine Einteilung entspricht \SI{5}{\milli\second}. Also ist eine Periode $4\cdot \SI{5}{\milli\second} = \SI{20}{\milli\second}$, also

  $$ f = \frac{1}{T}  = \frac{1}{\SI{0.02}{\second}} = \SI{50}{\hertz}$$
\end{sol}

\begin{sol}{EB411}
    Die Werte sind im Vergleich zur Frage \qref{EB409} nur im Komma verschoben:
    Eine Periode hat 4 Einteilungen und eine Einteilung entspricht \SI{0,03}{\micro\second}. Also ist eine Periode $4\cdot \SI{0,03}{\micro\second} = \SI{0,012}{\micro\second}$, also

  $$ f = \frac{1}{T}  = \frac{1}{\SI{0.00000012}{\second}} = \SI{8,33}{\mega\hertz}$$
\end{sol}





\begin{sol}{EI301}    
    Eine Periode hat 8 Einteilungen und eine Einteilung entspricht \SI{0,5}{\milli\second}. Also ist eine Periode $8\cdot \SI{0,5}{\milli\second} = \SI{4}{\milli\second}$.
\end{sol}


\begin{sol}{EI302}    
    Eine Periode hat 8 Einteilungen und eine Einteilung entspricht \SI{0,5}{\milli\second}. Also ist eine Periode $8\cdot \SI{0,5}{\milli\second} = \SI{4}{\milli\second}$.
 $$ f = \frac{1}{T}  = \frac{1}{\SI{0.004}{\second}} = \SI{250}{\hertz}$$
\end{sol}


\begin{sol}{EI303}    
    Wenn wir in etwa die Mitte des Signalausschlags betrachten, so hat er eine Länge von 4 Einteilungen.
   Eine Einteilung entspricht \SI{50}{\milli\second}. Also $4\cdot \SI{50}{\milli\second} = \SI{200}{\milli\second}$.
 \end{sol}

 \begin{sol}{EI304}    
    Für NF-Verzerrungen reicht sicherlich bereits ein einfaches Oszilloskop.
 \end{sol}

\begin{ohmchapter}
    In der Formelsammlung findest Du im Abschnitt \qq{Wechselspannung} die folgenden Formeln zur Umrechnung von Frequenz und Periodendauer:

    \begin{description}
    \item[Periodendauer] $$ T=\frac{1}{f} $$
    \item[Frequenz] $$ f = \frac{1}{T} $$
    \end{description}
\end{ohmchapter}   

\subsection{SMD-Widerstände}

 \begin{sol}{EC115}    
    Rechnung: $ 10 \cdot 10^3 = \SI{10}{\kilo\ohm}$
 \end{sol}

  \begin{sol}{EC116}    
    Rechnung: $ 22 \cdot 10^1 = \SI{220}{\ohm}$
 \end{sol}

  \begin{sol}{EC116}    
    Rechnung: $ 22 \cdot 10^1 = \SI{220}{\ohm}$
 \end{sol}

  \begin{sol}{EC117}    
    Rechnung: $ 22 \cdot 10^3 = \SI{22}{\kilo\ohm}$
 \end{sol}

\begin{ohmchapter}
Auf SMD-Widerständen ist der Wert in Form von Zahlen abgedruckt, wobei die letzte Ziffer die Zehnerpotenz angibt.
\end{ohmchapter}   

\subsection{Widerstandsmaterialien}

  \begin{sol}{EC101}    
    \qq{Hochlastwiderstände}: Drahtwiderstände.
  \end{sol}

  \begin{sol}{EC102}    
    \qq{Präzisionswiderstände}: Metallschichtwiderstände.
  \end{sol}

  \begin{sol}{EC103}    
    \qq{Induktionsarm + > \SI{30}{\mega\hertz}}: Metalloxidschichtwiderständen.
  \end{sol}

  \begin{sol}{EC103}    
    \qq{Induktionsarm + > \SI{30}{\mega\hertz}}: Metalloxidschichtwiderständen.
  \end{sol}

  \begin{sol}{EC104}    
    Im VHF- und UHF-Bereich ist es kritisch, dass Bauteile nur eine geringe Eigeninduktivität und Eigenkapazität haben, weil sie sonst als echte Antenne arbeiten würden.
  \end{sol}

   \begin{sol}{EC105}    
      Kohleschichtwiderstände sind besonders gut für hohe Frequenzen geeignet. 10 parallel geschaltete \SI{500}{\ohm} Widerstände ergeben \SI{50}{\ohm} Impedanz.
   \end{sol}

   \begin{sol}{EC106}    
      Analog zu Frage \qref{EC105} verwenden wir Kohleschichtwiderstände. In Parallelschaltung verteilt sich die Leistung auf die Widerstände, dies ist also besser für eine Dummy-Load.
   \end{sol}

  \begin{sol}{EC107}    
    \qq{Beste Eigenschaften + VHF}: Metalloxidschichtwiderständen (ungewendelt: Induktionsarm)
  \end{sol}

  \begin{sol}{EC108}    
    \qq{Temperaturmessung}: NTC-Widerstände.
  \end{sol}

\begin{ohmchapter}
\begin{description}
    \item[NTC / PTC] Heißleiter (NTC – Negative Temperature Coefficient) und Kaltleiter (PTC – Positive Temperature Coefficient) sind Widerstände, deren elektrischer Widerstand sich bei Temperaturänderungen signifikant ändert. Der Hauptunterschied liegt in der Richtung dieser Änderung.
    \begin{description}
         \item[NTC] Temperatur steigt: Widerstand sinkt.
         \item[PTC] Temperatur steigt: Widerstand steigt. 
    \end{description}  
    \item[Drahtwiderstände] Drahtwiderstände werden dazu z.B. aus Manganin und Konstantan hergestellt. Sie eignen sich als \textbf{Hochlastwiderstände}.
    \item[Metallschichtwiderstände] Metallschichtwiderstände können mit hoher Genauigkeit, d.h. mit geringer Fertigungstoleranz, gefertigt werden. Sie eignen sich als \textbf{Präzisionswiderstände}.
    \item[Kohleschicht- und Metalloxidschichtwiderständen] Bei Kohleschicht- und Metalloxidschichtwiderständen wird das Widerstandsmaterial in dünnen Schichten auf ein Trägermaterial aufgebracht. Beide Widerstandsarten sind \textbf{induktionsarm} und temperaturunabhängig und eignen sich besonders für den Einsatz bei Frequenzen oberhalb von \SI{30}{\mega\hertz}.
\end{description}   

\end{ohmchapter}  

\subsection{Widerstandstoleranzen}

  \begin{sol}{EC112}    
    Wir rechnen $\SI{5600}{\ohm} \cdot 1.1 = \SI{6160}{\ohm}$.    
  \end{sol}

  \begin{sol}{EC113}   
    Der silberne Farbring entspricht einer Toleranz von 10\%. Deshalb ist die Rechnung identisch zur Frage \qref{EC112}.
    Wir rechnen $\SI{5600}{\ohm} \cdot 1.1 = \SI{6160}{\ohm}$.    
  \end{sol}  
\begin{ohmchapter}
\end{ohmchapter}  

\subsection{Heißleiter und Kaltleiter}

  \begin{sol}{EC109}    
    Pfeile nach oben und unten: NTC.
  \end{sol}

  \begin{sol}{EC110}    
    NTC: Schaltzeichen mit Pfeil nach oben und unten.
  \end{sol}

  \begin{sol}{EC111}    
    PTC: Schaltzeichen mit zwei Pfeilen nach oben.
  \end{sol}  

\begin{ohmchapter}
\begin{description}
    \item[NTC / PTC] Heißleiter (NTC – Negative Temperature Coefficient) und Kaltleiter (PTC – Positive Temperature Coefficient) sind Widerstände, deren elektrischer Widerstand sich bei Temperaturänderungen signifikant ändert. Der Hauptunterschied liegt in der Richtung dieser Änderung.
    \begin{description}
         \item[NTC] Temperatur steigt: Widerstand sinkt. Das Schaltzeichen hat einen \textbf{Pfeil nach oben und einen nach unten}.
         \item[PTC] Temperatur steigt: Widerstand steigt. Das Schaltzeichen zeigt \textbf{zwei Pfeile nach oben}.
    \end{description}  
\end{description}      
\end{ohmchapter}  

\subsection{Leistung II}

  \begin{sol}{EB501}    
    Die Abkürzung PEP steht für den englischen Ausdruck \qq{peak envelope power}.
    Es geht auf Deutsch übersetzt um die \textbf{Modulationshüllkurve}. Du solltest Dir auch merken, dass PEP immer direkt am Ausgang des Funkgeräts bestimmt wird (d.h. ohne Berücksichtigung von Dämpfung oder Antennengewinn). Gemessen wird an einem sogenannten \textbf{reellen Abschlusswiderstand}. Du kannst Dir unter reellem Abschlusswiderstand vereinfacht eine Dummyload vorstellen. Tatsächlich ist dies im Gegensatz zu einem Blindwiderstand definiert, diesen werden wir aber erst in der Klasse A genau definieren. Zum Beantworten der Frage habe ich Dir Schlüsselworte markiert.
  \end{sol}  


    \begin{sol}{EB502}    
 Bei der \qq{mittlere(n) Leistung} geht es, wie der Name bereits vermuten lässt, um eine \textbf{durchschnittliche Leistung}. D. h., es wird über einen Zeitraum gemittelt und der Zeitraum sollte ausreichend lang sein. In der Musterantwort soll sie \textbf{im Verhältnis zur Periode der tiefsten Modulationsfrequenz} ausreichend lang sein.
    \end{sol}  

   \begin{sol}{EB503}    
    Wir hatten uns damit bereits im Kapitel zur Wechselspannung beschäftigt. Dies ist auch einer der Gründe, warum die effektive Spannung $U_\text{eff}$ für uns wichtig war.
   \end{sol}  


   \begin{sol}{EB504}    
    Die richtige Antwort $U=\sqrt{P\cdot R}$ steht in der Formelsammlung im Abschnitt \qq{Leistung}.
   \end{sol}  

   \begin{sol}{EB505}    
    Die richtigen Formeln $I=\frac{P}{R}$ und $U=\sqrt{P\cdot R}$ stehen in der Formelsammlung im Abschnitt \qq{Leistung}.
   \end{sol}  

   \begin{sol}{EB506}    
    Die richtigen Formeln stehen so nicht in der Formelsammlung, da hier nicht nach $R$ aufgelöst wurde.
    Dies lässt sich aber mit den Formeln in der Formelsammlung in Abschnitt \qq{Leistung} schnell erledigen:
    $$ P = \frac{U^2}{R} \Rightarrow R = \frac{U^2}{P} $$
    $$ P = I^2 \cdot R \Rightarrow R = \frac{P}{I^2} $$
   \end{sol}  


\begin{sol}{EB507}  
$$P = \frac{U^2}{R} = \frac{\SI{100}{\volt}^2}{\SI{50}{\ohm}} = \SI{200}{\watt}$$
\end{sol}  

\begin{sol}{EB508}  
$$P = I^2\cdot R= \SI{2}{\ampere}^2\cdot \SI{50}{\ohm} = \SI{200}{\watt}$$
\end{sol}  

\begin{sol}{EB509}  
$$P = \frac{U^2}{R} = \frac{\SI{10}{\volt}^2}{\SI{100}{\ohm}}= \SI{1}{\watt}$$
\end{sol} 


\begin{sol}{EB510}  
$$U = \sqrt{P\cdot R} = \sqrt{\SI{1}{\watt}\cdot \SI{10000}{\ohm}}  =\SI{100}{\volt} $$
\end{sol} 


\begin{sol}{EB511}  
$$U = \sqrt{P\cdot R} = \sqrt{\SI{6}{\watt}\cdot \SI{100000}{\ohm}}  \approx \SI{775}{\volt}$$
\end{sol} 

\begin{sol}{EB512}      
$$I = \sqrt{\frac{P}{R}} = \sqrt{\frac{\SI{23}{\watt}}{\SI{120}{\ohm}}}  = \SI{438}{\milli\ampere}$$
\end{sol} 


\begin{sol}{EB513}  
    $$\hat{U} = \SI{25}{\volt} / 2 = \SI{12.5}{\volt} $$
    $$U_\text{eff} = \frac{\hat{U}}{\sqrt{2}} = \frac{\SI{12.5}{\volt}}{\sqrt{2}}  = \SI{8.8}{\volt}$$
$$I = \frac{U}{R} = \frac{ \SI{8.8}{\volt}}{\SI{1000}{\ohm}}= \SI{8.8}{\milli\ampere}$$
\end{sol} 

\begin{sol}{EB514}  
 Jeder der 11 Widerstände ist mit jeweils \SI{5}{\watt} belastbar. Also $11\cdot \SI{5}{\watt} = \SI{55}{\watt}$
\end{sol} 

\begin{ohmchapter}    
\end{ohmchapter}  

\subsection{Dezibel I}

  \begin{sol}{EA107}    
    Leistungsverdopplung kommt immer wieder vor, also gut auswendig zu wissen, dass dies \SI{3}{\decibel} entspricht. Falls Du es Dir nicht gemerkt hast, findest Du den Wert in der Formelsammlung im Abschnitt \qq{Pegel}.
 \end{sol}
\begin{ohmchapter}
\end{ohmchapter}  